\section{PORS+FP}
PORS (PRNG to Obtain a Random Subset) is a few-time signature scheme that uses a \emph{single} Merkle tree over \texttt{t}
leaves. PORS reveals \texttt{k} leaves from a single tree and uses the Octopus algorithm~\cite{pors_octopus2018} to compute the minimal
set of authentication nodes needed to verify all \texttt{k} leaves simultaneously,
exploiting shared nodes to shorten the proof. PORS+FP~\cite{pors_fp2025} adds
\emph{forced pruning} on top: the signer grinds a counter until the resulting
$k$-subset of leaves produces an Octopus authentication set of size at most
\texttt{m\_max} nodes, bounding the signature to a fixed maximum size.

\subsection{PORS Structure}
 
PORS consists of a single Merkle tree over \texttt{t} leaves:
\begin{itemize}
    \item The tree has \texttt{t} leaves (not necessarily a power of~2; the tree is
          left-filled as described in Section~2.1)
    \item Each leaf is derived from \texttt{SK.seed} using \texttt{PRF} with address
          type \texttt{PORS\_PRF}
    \item The message together with a grinding randomness \texttt{R},
          is hashed to a n-byte digest which is then expanded into \texttt{k} distinct
          indices in $[t] = \{0,1,...,t-1\}$
    \item The Octopus algorithm computes the minimal authentication set \texttt{A} for those
          \texttt{k} leaves
\end{itemize}
 
The PORS+FP public key is the tweakable hash of the single tree root:
\begin{center}
    \texttt{PORS+FP PK = Tw(PK.seed, ADRS, root)},
\end{center}
using address type \texttt{PORS\_PK}. Forced pruning guarantees
$|A| \le m_{\max}$, bounding the signature to a fixed maximum size.

\subsection{Grinding Condition}
\label{subsec:pors_grind_condition}
 
PORS+FP uses the same two-call message-hash interface (Section~5.4.2):
\begin{verbatim}
    R      ← PRF_msg_ctr(SK.prf, PK.seed, opt_rand, ctr, m)
    digest ← H_msg_pors(ADRS, R, PK.seed, PK.root, m)
    A      ← pors_octopus(pors_digest_to_indices(digest))
    condition: |A| <= m_max
\end{verbatim}
 
The signer searches for an \texttt{R} that satisfies this condition by
iterating a counter inside \texttt{PRF\_msg\_ctr}. Each counter value produces
a fresh \texttt{R}, which is then evaluated via \texttt{H\_msg\_pors}. The
verifier only sees the final \texttt{R} that satisfied the condition and
recomputes the same digest deterministically. Expected grinding cost depends on the distribution of $|A|$; 

\subsection{Digest to Indices}
\label{subsec:pors_digest_to_indices}
 
The \texttt{pors\_digest\_to\_indices} function converts a message digest and a grinding
counter into \texttt{k} \emph{distinct} indices in $\{0,1,...,t-1\}$. Distinctness is required by the Octopus
algorithm. We use rejection sampling over SHA256 XOF output.
 
\begin{verbatim}
    Algorithm: pors_digest_to_indices(digest)
    Input:
        digest: 32-byte value
    Output:
        indices: array of k distinct integers in [0, t-1]
        xof_block: XOF output with size 256*blk

    1. b       ← roundup(log2(t))         // bits per candidate
    2. c       ← floor(256 / b)        // number of candidates per 256-bit block
    3. total_bits ← (roundup(2^b * k / t) + margin) * b + offset + hsl // XOF output length
    4. min_blk ← roundup(total_bits/256) // expected number of XOF blocks
    5. indices ← []
    6. xof_out ← ""
    7. setTypeAndClear(ADRS, PORS_XOF)
    8. for blk from 0 to 2^32 - 1:
        a. block ← H_xof_pors(ADRS, PKseed, digest, blk)
        b. xof_out ← xof_out || block
        c. for i from 0 to c - 1:
            if |indices| == k:
                break
            candidate ← extract_bits(block, i * b, b)
            if candidate < t and candidate not in indices:
                indices ← indices || [candidate]
        d. if |indices| == k and blk = min_blk - 1:
            return (indices, xof_out[0:total_bits])
    9. abort with error
\end{verbatim}

The length of the XOF output is \texttt{total\_bits} (see Section~5.4.2). 
In our case, for chosen parameters of SHRINCS-B and SHRINCS-L, and \texttt{margin = 7}, \texttt{offset = 16}, we get the follwing length: \texttt{total\_bits = (roundup(2\string^24 * 6 / 9245141) + 7) * 24 + 16 + 24 = 472} bits or 59 bytes.
For SHRINCS-B32 parameters we set \texttt{margin = 10}, \texttt{offset = 16}, and in the result is: \texttt{total\_bits = (roundup(2\string^17 * 11 / 109571) + 10) * 17 + 16 + 32 = 456} bits or 57 bytes.

\paragraph{Note:} \texttt{extract\_bits(x, offset, length)} extracts length bits from byte string \texttt{x} starting at bit position offset, returned as an unsigned integer. Bits are numbered in big-endian order: bit \texttt{0} is the most significant bit of byte \texttt{x[0]}.

\subsection{Octopus for the PORS+FP Tree}
 
The Octopus algorithm~\cite{pors_octopus2018} computes the minimal set of authentication
nodes whose values are sufficient to recompute the tree root from \texttt{k} revealed
leaves. It is invoked during grinding (to check the size condition), during signing
(to determine which node values to include in the signature), and implicitly during
verification (to determine which nodes to read from the signature). Because all
three parties run the same algorithm on the same index set, the authentication node
ordering is unambiguous.
 
For an imperfect, left-filled tree of \texttt{t} leaves we initialise the algorithm by
partitioning the \texttt{k} input indices by their actual depth in the tree,
then proceed identically to the standard Octopus formulation.
 
\begin{verbatim}
    Algorithm: pors_octopus(indices)
    Input:
        indices: array of k distinct leaf indices in [0, t-1], sorted ascending
    Output:
        A:       array of authentication node (lvl, idx) pairs

    1. h ← ceil(log2(t))
    2. s ← t - 2^(h-1)
    3. I ← []
    4. P ← []
    5. A ← []
    6. for each i in indices:
        a. if i < 2 * s:
            I ← I || [(0, i)]
        b. else:
            P ← P || [(1, i - s)]          
    7. for current_lvl from 0 to h - 1:
        a. next_P ← []
        b. i ← 0
        c. while i < |I|:
            (lvl, idx) ← I[i]
            sib_idx ← idx XOR 1
            if i + 1 < |I| and I[i+1][1] == sib_idx:
                i ← i + 2
            else:
                A ← A || [(current_lvl, sib_idx)]
                i ← i + 1
            next_P ← next_P || [(current_lvl + 1, idx >> 1)]
            
        d. I ← next_P || P
        e. P ← []
        f. if |I| == 1 and I[0][1] == 0:
            break   
    8. return A
\end{verbatim}
 
\paragraph{Note:} When \texttt{t} is a power of~2, $s = t/2$ and all leaves are at depth
\texttt{h}, reducing step~4 to $I = \{(h, i) : i \in \texttt{indices}\}$ and recovering
the standard Octopus formulation. The algorithm performs no hash evaluations; its
cost is $O(k \log t)$ set operations.

\subsection{Grinding}
 
The \texttt{pors\_grind} function searches for a randomness \texttt{R} such that
the Octopus authentication set for the resulting index set has size at most
\texttt{m\_max}. The winning \texttt{R} is returned and included directly in the enclosing 
SPHINCS signature; the verifier recomputes \texttt{digest} and \texttt{indices} 
deterministically from \texttt{R} and \texttt{m}.
 
\begin{verbatim}
    Algorithm: pors_grind(m, SK.prf, PK.seed, PK.root, ADRS, opt_rand)
    Input:
        m:       message
        SK.prf:  PRF key for randomness generation
        PK.seed: public seed
        PK.root: public key root
        ADRS:    tweak
        opt_rand: optional randomness (n bytes)
    Output:
        indices: array of k integers in [0, t-1]
        R:       randomness satisfying the forced-pruning condition
        digest:  H_msg_pors output for that R
 
    1. setTypeAndClear(ADRS, SL_H_MSG)
    2. for ctr from 0 to 2^32 - 1:
        a. R      ← PRF_msg_ctr(SK.prf, PK.seed, opt_rand, ctr, m)
        b. digest ← H_msg_pors(ADRS, R, PK.seed, PK.root, m)
        c. indices ← pors_digest_to_indices(digest)
        d. A       ← pors_octopus(indices)
        e. if |A| <= m_max:
               return (indices, R, digest)
    3. abort with error
\end{verbatim}
 
\paragraph{Note:} Each iteration costs one \texttt{PRF\_msg\_ctr} evaluation plus
one \texttt{H\_msg\_pors} evaluation and dominates the Octopus authentication cost. The address type \texttt{SL\_H\_MSG} is set
at the start and remains fixed for the duration of grinding.

\subsection{Secret Key Generation}
 
The \texttt{pors\_SKgen} function derives the secret value for a single PORS leaf.
 
\begin{verbatim}
    Algorithm: pors_SKgen(SK.seed, PK.seed, ADRS, leaf_idx)
    Input:
        SK.seed:  secret seed
        PK.seed:  public seed
        ADRS:     tweak
        leaf_idx: leaf index in [0, t-1]
    Output:
        sk: secret value (n bytes)
 
    1. setTypeAndClear(ADRS, PORS_PRF)
    2. setKeyPairAddress(ADRS, 0)
    3. setTreeIndex(ADRS, leaf_idx)
    4. sk ← PRF(PK.seed, ADRS, SK.seed)
    5. return sk
\end{verbatim}

\subsection{TreeHash}

The \texttt{pors\_treehash} function recursively computes the value associated with
node \texttt{(target\_height, idx)} in the single PORS+FP Merkle tree, where
\texttt{target\_height} is the height of the node counted upward from the leaves and \texttt{idx}
is its horizontal index at that height. The tree is left-filled with \texttt{t} leaves: letting $h = \lceil \log t \rceil$ and $s = t - 2^{h-1}$, the
leftmost $2s$ leaves sit at depth $h$ (or height~0) and the
remaining $t - 2s$ leaves sit at depth $h - 1$ (height~1).

\begin{verbatim}
    Algorithm: pors_treehash(SK.seed, PK.seed, ADRS, target_height, idx)
    Input:
        SK.seed, PK.seed: seeds
        ADRS:             tweak
        target_height:    height of output node
        idx:        horizontal index of the node at target_height
    Output:
        node: hash value (n bytes)

    1. h ← ceil(log2(t))
    2. s ← t - 2^(h-1)
    3. if target_height == 0:
        a. sk ← pors_SKgen(SK.seed, PK.seed, ADRS, idx)
        b. setTypeAndClear(ADRS, PORS_HASH)
        c. setKeyPairAddress(ADRS, 0)
        d. setTreeHeight(ADRS, 0)
        e. setTreeIndex(ADRS, idx)
        f. return Tw(PK.seed, ADRS, sk)
    4. if target_height == 1 and idx >= s:
        a. leaf_idx ← s + idx
        b. sk ← pors_SKgen(SK.seed, PK.seed, ADRS, idx)
        c. setTypeAndClear(ADRS, PORS_HASH)
        d. setKeyPairAddress(ADRS, 0)
        e. setTreeHeight(ADRS, 0)
        f. setTreeIndex(ADRS, idx)
        g. return Tw(PK.seed, ADRS, sk)
    5. left  ← pors_treehash_paper(SK.seed, PK.seed, ADRS, target_height - 1, 2*idx)
    6. right ← pors_treehash_paper(SK.seed, PK.seed, ADRS, target_height - 1, 2*idx + 1)
    7. setTypeAndClear(ADRS, PORS_TREE)
    8. setKeyPairAddress(ADRS, 0)
    9. setTreeHeight(ADRS, target_height)
    10. setTreeIndex(ADRS, idx)
    11. return Tw(PK.seed, ADRS, left || right)
\end{verbatim}

\subsection{PORS+FP Authentication path}

The \texttt{pors\_auth\_path} function generates the minimal set of authentication nodes required to verify all \texttt{k} revealed leaves simultaneously within the single PORS+FP Merkle tree. 
Rather than authenticating each leaf independently, it employs the Octopus algorithm to identify shared internal nodes, which significantly shortens the resulting proof. 
For every node coordinate \texttt{(lvl, idx)} returned by the Octopus algorithm, the function computes the required hash value by invoking \texttt{pors\_treehash}. 
In this context, \texttt{lvl} specifies the height of the authentication node, and \texttt{idx} represents the leftmost leaf index of the subtree rooted at that node.
\begin{verbatim}
    Algorithm: pors_auth_path(SK.seed, PK.seed, ADRS, indices)
        Input:
            SK.seed, PK.seed: seeds
            ADRS:    tweak
            indices: array of k distinct leaf indices in [0, t-1], sorted ascending
        Output:
            auth:    authentication path (sequence of n-byte hashes)

        1. auth ← []
        2. A ← pors_octopus(indices)
        3. for each (lvl, idx) in A:  
            a. node ← pors_treehash(SK.seed, PK.seed, ADRS, lvl, idx)
            b. auth ← auth || node
        4. return auth
\end{verbatim}

\subsection{Signature generation}
The \texttt{pors\_sign} function produces a PORS+FP signature. Due to forced pruning, the Octopus authentication set is guaranteed to contain at most $m_{\max}$ nodes. The signature format is:
\begin{center}
    \texttt{R || sk\_\{i\_1\} || ... || sk\_\{i\_k\} || y\_\{a\_1\} || ... || y\_\{a\_{|A|}\}}
\end{center}
where $\{i_1,\dots,i_k\} = I$ are the \texttt{k} revealed leaf indices and $\{a_1,\dots,a_{|A|}\} = A$ are the Octopus authentication node labels (in the canonical Octopus traversal order).
We pad additional n-byte 0 strings to make the length of the signature be fixed with exactly size of \texttt{(k + m\_max) * n} bytes.

\begin{verbatim}
    Algorithm: pors_sign(m, SK.seed, SK.prf, PK.seed, PK.root, ADRS)
    Input:
        m: message
        SK.seed: secret seed
        SK.prf: PRF key
        PK.seed: public seed
        PK.root: public key root
        ADRS: tweak
    Output:
        sig: PORS+FP signature
        digest: the digest that satisfied the grinding condition

    1. opt_rand ← random(n) or PK.seed
    2. (indices, R, digest) ← pors_grind(m, SK.prf, PK.seed, PK.root, ADRS, opt_rand)
    3. sig ← R
    4. for each i in indices (in increasing order):
        a. sk_i ← pors_SKgen(SK.seed, PK.seed, ADRS, i)
        b. sig  ← sig || sk_i
    5. auth_path ← pors_auth_path(SK.seed, PK.seed, ADRS, indices)
    6. sig ← sig || auth_path
    7. for ctr from |auth_path| to m_max - 1:
        a. sig ← sig || toByte(0, n)
    8. return (sig, digest)
\end{verbatim} 

\subsection{Public Key Recovery}
The \texttt{pors\_pkFromSig} function recovers the PORS+FP public key from a signature. 
For each of the \texttt{k} revealed leaves, it hashes the secret key to obtain the leaf value, then runs the Octopus-guided root recomputation using the supplied authentication nodes. 
Verification succeeds if and only if the recovered root matches \texttt{PK.root}.

\begin{verbatim}
   Algorithm: pors_pkFromSig(sig, indices, PK.seed, PK.root, ADRS)
    Input:
        sig:     PORS+FP signature (R || secret leaves || auth nodes)
        indices: PORS+FP k indices from digest
        PK.seed: public seed
        PK.root: public key root
        ADRS:    tweak (layer and tree_address set by caller)   
    Output:
        pk: PORS public key

    1. offset ← R_SIZE
    2. h ← ceil(log2(t))
    3. s ← t - 2^(h-1)
    4. I ← []
    5. P ← []
    6. for i from 0 to k - 1:
            a. sk_i ← sig[offset : offset + n]
            b. offset ← offset + n
            c. setTypeAndClear(ADRS, PORS_HASH)
            d. setKeyPairAddress(ADRS, 0)
            e. setTreeHeight(ADRS, 0)
            f. setTreeIndex(ADRS, indices[i])
            g. val ← Tw(PK.seed, ADRS, sk_i)
            h. if indices[i] < 2*s:
                I ← I || [(0, indices[i],     val)]
            i. else:
                P ← P || [(1, indices[i] - s, val)]
    7. for current_lvl from 0 to h - 1:
        a. paired ← [false] * |I|
        b. for i from 0 to |I| - 1:
            if i + 1 < |I| and I[i+1][1] == I[i][1] XOR 1:
                paired[i]     ← true
                paired[i + 1] ← true
        c. next_P ← []
        d. for i from 0 to |I| - 1:
            (lvl, idx, val) ← I[i]
            if paired[i] and (idx AND 1) == 0:
                continue                 
            setTypeAndClear(ADRS, PORS_TREE)
            setKeyPairAddress(ADRS, 0)
            setTreeHeight(ADRS, current_lvl + 1)
            setTreeIndex(ADRS, idx >> 1)    
            if paired[i]: 
                sib_val ← I[i-1][2]
                parent_val ← Tw(PK.seed, ADRS, sib_val || val)
            else:
                auth_val ← sig[offset : offset + n]
                offset ← offset + n
                if (idx AND 1) == 0:
                    parent_val ← Tw(PK.seed, ADRS, val || auth_val)
                else:
                    parent_val ← Tw(PK.seed, ADRS, auth_val || val)
            
            next_P ← next_P || [(current_lvl + 1, idx >> 1, parent_val)]

        e. I ← next_P || P
        f. P ← []
    8. root ← I[0][2]
    9. setTypeAndClear(ADRS, PORS_PK)
    10. pk ← Tw(PK.seed, ADRS, root)
    11. return pk
\end{verbatim}